\chapter{Performance}\label{ch:performance}
\sloppy

Dylan was designed to be both dynamic and efficient, and the compiler to
C aims at its original targets: small and fast. This chapter reports how
fast compiled programs are on a set of benchmarks, explains what makes a
Dylan program fast or slow under each implementation, and shows how to
measure a program and tune the garbage collector. Most of the advice
follows from how the compiler works (Chapter~\ref{ch:internals}); the
measurements here show how much each point matters.

\begin{notebox}
Every time in this chapter was measured on one machine: an Intel Core
i9-14900KF running Fedora Linux, with SBCL~2.6.3 and GCC~16.2.1. Times
depend on the processor, the C compiler, the collector's version and the
load on the machine, and vary by several percent from run to run.
Treat them as an indication of proportions, not as absolute figures, and
run the benchmarks on your own machine for numbers that apply to you.
\end{notebox}

\section{The benchmarks}\label{sec:perf-benchmarks}

\index{benchmarks}
The directory \file{bench/} holds eleven small programs, each written
three times: in Dylan (\file{benchmarks.dyl}), in plain Common Lisp
(\file{reference.lisp}) and in plain C (\file{reference.c}). Each answers
a value, and the three versions must agree.

\begin{center}
\small
\begin{tabular}{>{\raggedright\arraybackslash}p{0.17\textwidth}%
                >{\raggedright\arraybackslash}p{0.72\textwidth}}
\toprule
Name & What it does \\
\midrule
\code{table} & insert into and look up in a \code{<table>} \\
\code{sort} & sort a list of pseudo-random integers \\
\code{string} & build and compare strings \\
\code{deque} & push and pop at both ends of a \code{<deque>} \\
\code{stretchy} & grow a \code{<stretchy-vector>} and walk it \\
\code{range} & iterate over ranges \\
\code{exit} & leave a computation early with \code{bind-exit} \\
\code{condition} & signal errors and handle them with \code{handler-case} \\
\code{next-method} & call down a chain of methods with \code{next-method} \\
\code{tak} & the Takeuchi function: integer arithmetic and calls \\
\code{float} & points of a grid inside the Mandelbrot set: floating
  point \\
\bottomrule
\end{tabular}
\end{center}

Two scripts run them. Each prints, for every benchmark, the best of
three runs.

\begin{shell}
$ sbcl --script bench/run-c.lisp          # compiled Dylan vs C and Lisp
$ sbcl --script bench/run.lisp            # interactive Dylan vs Lisp
$ sbcl --script bench/run.lisp table tak  # just these
\end{shell}

\noindent Their builds go in \file{bench/build/}. A line marked
\code{MISMATCH} means the versions answered differently, which would be
a bug.

\index{BENCH\_SCALE@\code{BENCH\_SCALE}}
At their usual sizes several benchmarks finish in a few milliseconds,
too quickly to compare reliably. The environment variable
\code{BENCH\_SCALE}, if set to an integer, multiplies every size, in
both scripts and in \file{reference.c}:

\begin{shell}
$ BENCH_SCALE=10 sbcl --script bench/run-c.lisp
\end{shell}

\noindent The file \file{SpeedComparison.md} at the top of the
repository puts the three ways of running the benchmarks --- plain SBCL,
the Common Lisp implementation and the compiler to C --- side by side,
with hand-written C, at \code{BENCH\_SCALE=10}, and discusses the
results. The tables below were run at the usual sizes.

\subsection{The compiler to C}

\cmd{bench/run-c.lisp} compiles the Dylan with \cmd{dylanc} (at
\cmd{-O2}), compiles \file{reference.c} with \cmd{cc -O2}, and loads
\file{reference.lisp} into SBCL. A run gave:

\begin{shell}
benchmark     Dylan->C ms     C ms      D/C    Lisp ms   D/Lisp
table                  13        6     2.17        8.0     1.63
sort                   29       13     2.23       26.0     1.12
string                 11        2     5.50        7.0     1.57
deque                   6        2     3.00        2.0     3.00
stretchy                3        0        -        2.0     1.50
range                   5        0        -        2.0     2.50
exit                    1        3     0.33        1.0     1.00
condition               4        0        -        1.0     4.00
next-method             5        2     2.50        8.0     0.62
tak                    22       17     1.29       35.0     0.63
float                   3        2     1.50       55.0     0.05
\end{shell}

The columns are the compiled Dylan's time, the C's, their ratio, plain
SBCL's time and the ratio to it. A ratio is shown as \code{-} when the
other time is under a millisecond and the ratio would mean nothing.

Read it this way. On calls and integer arithmetic (\code{tak}) and on
floating point (\code{float}), compiled Dylan takes 1.3 to 1.5 times as
long as hand-written C. On the collections it is two to three times the C's
time, and on strings about five times. Non-local exits (\code{exit}) are
\emph{faster} than the C, because the Dylan exit becomes a jump where
the C uses \code{setjmp}. Against plain SBCL, compiled Dylan is faster
on calls, \code{next-method} and floats, and slower on the collection
benchmarks by up to three times. The very short benchmarks
(\code{stretchy}, \code{range}, \code{condition}) take a few
milliseconds where the C takes under one, so their ratios vary most from
run to run; in other runs the \code{string} ratio was 6.5 and the
\code{table} ratio 1.6.

\subsection{The interactive implementation}

\cmd{bench/run.lisp} compiles \file{benchmarks.dyl} with the interactive
implementation's \code{dylan:compile-file}, compiles the resulting Lisp
with SBCL, and compares it with \file{reference.lisp}:

\begin{shell}
benchmark         size      dylan       lisp    ratio
table           200000     0.076s     0.007s    10.9x
sort            100000     0.083s     0.025s     3.3x
string          200000     0.050s     0.005s    10.0x
deque           300000     0.029s     0.002s    14.5x
stretchy        300000     0.026s     0.002s    13.0x
range          1000000     0.048s     0.003s    16.0x
exit            300000     0.027s     0.001s    27.0x
condition        30000     0.010s     0.000s    10.0x
next-method     300000     0.039s     0.008s     4.9x
tak                300     0.469s     0.036s    13.0x
float              300     0.089s     0.054s     1.6x
\end{shell}

The interactive implementation is typically ten to fifteen times slower
than plain Lisp, and so ten to twenty times slower than the compiler to
C. That is the price of a language whose classes and methods can change
while it runs: every call of a generic function must allow for methods
added since the last call. The interactive implementation is for
developing and exploring; the compiler is for running.

\section{Writing fast Dylan for the compiler}\label{sec:perf-c}

The compiler is fastest where it knows the classes of the values it
works with. Then it calls methods directly, does arithmetic and
collection access inline without tests, and keeps floats unboxed. Each
point below says how to give it that knowledge, with a measurement. The
programs are small, and each was timed as a whole, run as a separate
executable built at \cmd{-O2}.

\subsection{Prefer \code{define-method} to \code{define} of a \code{method}}

How a function is defined matters more than you might expect. These
four versions of \code{tak}, computing \code{(tak 30 22 12)}, differ
only in their first lines:

\begin{center}
\small
\begin{tabular}{lr}
\toprule
Definition & Time \\
\midrule
\code{(define tak (method (x y z) ...))} & 0.04 s \\
\code{(define-method tak (x y z) ...)} & 0.04 s \\
\code{(define-method tak ((x <integer>) (y <integer>) (z <integer>)) ...)}
  & 0.08 s \\
\code{(define tak (method ((x <integer>) (y <integer>) (z <integer>)) ...))}
  & 0.30 s \\
\bottomrule
\end{tabular}
\end{center}

\begin{itemize}
\item An unspecialized \code{define-method} with only one method is not
  dispatched at all, and a \code{define} of a \code{method} with no
  specializers is called directly; both are as fast as the program can
  be.
\item Specializing the only method on \code{<integer>} makes every call
  go through the dispatch function, which tests the argument classes.
  Inside the method the arguments are then known to be integers, which
  helps a method that does much work with them, but here the method is
  tiny and the dispatch doubles its cost.
\item A \code{method} with specializers stored with \code{define} is the
  slowest: it is a closure that checks its own arguments on every call,
  is called through the variable, and does not pass on to the compiler
  the knowledge its specializers give.
\end{itemize}

So: specialize methods when you need dispatch, or when the method does
enough with its arguments to profit from knowing their classes; do not
specialize a lone method merely to document its argument types; and
define functions with \code{define-method}, or with \code{define} of an
unspecialized \code{method}.

\subsection{Let the compiler know classes: declarations in \code{bind}}

\index{bind@\code{bind}!type declarations}
A value fetched from a collection, or returned by a function the
compiler cannot see into, has no known class. Every operation on it is
then preceded by tests of its class, and a float result must be boxed
because the compiler cannot be sure it is a float. A type declaration in
\code{bind} (\dylanbook{chapter 4}) tells the compiler the class; it is
checked once, when the binding is made. Summing squares of the floats in
a vector of 1000 elements, 100\,000 times:

\begin{dylan}
(define-method sum-squares (v)
  (bind ((s 0.0))
    (for-each ((x v)) ()
      (set! s (+ s (* x x))))
    (> s 0.0)))                                   ; 1.06 s
\end{dylan}

\begin{dylan}
(define-method sum-squares (v)
  (bind ((s 0.0))
    (for-each ((x v)) ()
      (bind (((y <double-float>) x))
        (set! s (+ s (* y y)))))
    (> s 0.0)))                                   ; 0.20 s
\end{dylan}

In the first version \code{(* x x)} tests whether \code{x} is an integer
or a float and allocates a box for the product; in the second it is one
machine multiplication, and nothing is allocated. Specializing a method
on the class of its argument has the same effect inside the method; so
does a \code{type:} on a slot, for the values its getter answers.

Collections are no different. A \code{for-each} over a parameter of
unknown class tests, on each step, whether it is walking a list or a
vector; over one known to be a \code{<simple-object-vector>} it is a
plain indexed loop. Summing a vector of 1000 integers 100\,000 times took
0.16~s with an unspecialized parameter and 0.10~s with the parameter
specialized on \code{<simple-object-vector>}.

\subsection{Keep floats unboxed}

\index{unboxing}
A \code{<double-float>} is a heap object, so a float computation that
creates a new one on every step spends most of its time allocating. The
compiler avoids this by keeping a float variable as a C \code{double}
when \emph{every} use of it wants a \code{double}
(Section~\ref{sec:internals-types}). Returning the variable, storing it,
or passing it to a function the compiler cannot see into each need a
box, and force one on every step. Compare (100 million steps each):

\begin{dylan}
(define-method norm-sum (scale n)
  (for ((i 0 (+ i 1))
        (s 0.0 (+ s (* scale i))))
       ((= i n) s)))                    ; returns s itself: 0.87 s

(define-method norm-sum (scale n)
  (for ((i 0 (+ i 1))
        (s 0.0 (+ s (* scale i))))
       ((= i n) (> s 1.0)))             ; only compares s: 0.04 s
\end{dylan}

The difference is not the work but the boxes. A hot loop should keep its
float accumulators in variables used only arithmetically, and box the
result once, after the loop, by passing it to a function or returning it
from an enclosing one.

\paragraph{Use float contagion, not \code{as}.}
\index{float contagion}
Mixing an integer into float arithmetic converts it automatically
(\dylanbook{chapter 14}), and the compiler does the conversion inline.
An explicit \code{(as <double-float> i)} calls the generic function
\code{as} instead. In the loop above, writing \code{(* scale (as
<double-float> i))} instead of \code{(* scale i)} made it two and a half
times slower (2.2~s against 0.87~s).

\subsection{Typed slots}

A slot declared with \code{type: <double-float>} lets the compiler know
the class of what its getter answers, so arithmetic on it needs no test.
The value stored in the slot is still a boxed float, so the gain is
modest when the program writes the slot often: stepping a particle's
position 100 million times took 1.18~s with untyped slots and 1.05~s
with typed ones.

\subsection{Local exits are free; escaping ones are not}

\index{bind-exit@\code{bind-exit}!cost}
A \code{bind-exit} whose exit procedure is called only directly within
the function that binds it --- an early return from a loop --- compiles
to a jump. One whose exit is passed to another function needs a frame in
the dynamic chain, \code{setjmp}, \code{longjmp}, and a closure
(Section~\ref{sec:internals-exits}). Searching a ten-element vector 30
million times, returning early:

\begin{dylan}
(define-method find-first (v (x <integer>))
  (bind-exit (return)
    (for-each ((e v)) ()
      (when (= e x) (return e)))
    #f))                                              ; 0.08 s

(define-method check (e x return) (when (= e x) (return e)))
(define-method find-first (v (x <integer>))
  (bind-exit (return)
    (for-each ((e v)) ()
      (check e x return))
    #f))                                              ; 1.48 s
\end{dylan}

The second is nearly twenty times slower. Keep the calls of an exit
procedure inside the function that binds it where you can. The same
applies to \code{unwind-protect} (and so to \code{with-lock} and
\code{dynamic-bind}, which use it), whose cleanup always needs a frame,
and to \code{handler-case}: a handled error costs a few hundred
nanoseconds.

\subsection{First-class calls}

A call of a function held in a variable or passed as an argument goes
through the function's uniform entry, with no inlining and no
specialized entry. In a hot loop it costs more than a direct call or
inline code, but often less than you might fear: summing squares of a
vector with \code{(reduce (method (acc x) (+ acc (* x x))) 0 v)} took
0.19~s where the equivalent \code{for-each} loop took 0.16~s. It matters
most when the function called is small and could otherwise have been
inlined. \code{apply} with a long list, and \code{\#rest} lists that
escape, allocate.

\subsection{\code{\#rest} and keywords}

A \code{\#rest} list that is only walked --- not stored, returned, or
captured by a closure --- is built on the C stack, for up to eight
elements, and costs no allocation. Keyword arguments to a function the
compiler can see are matched when compiling, with no list built at all.
So these are cheap in the common cases; a \code{\#rest} list that is
stored or returned is allocated on the heap.

\subsection{What the static program gives}

Everything in a compiled program is effectively sealed: the compiler
knows every class and every method. That is what allows it to do
\code{+} on integers inline, or \code{element} on a vector: it can see
that no method but the library's could apply to such arguments. If your
program defines methods on the library's generic functions for the
library's own classes --- a method on \code{binary+} for two integers,
say, or on \code{element} for \code{<simple-object-vector>} --- the
inline code for those classes is dropped, and every such call in the
program goes through dispatch. (With a method on \code{binary+} for two
integers, \code{+} of two integers calls it, as it must; the inline
cases for floats remain.) Methods for your own classes, or for
combinations the library does not handle inline, do not have this
effect.

\section{Writing fast Dylan for the interactive implementation}
\label{sec:perf-lisp}

The interactive implementation compiles Dylan to Common Lisp, which SBCL
compiles to machine code. Generic functions dispatch through caches;
calls of known functions are made directly; arithmetic on numbers,
\code{element} on simple vectors and strings, and iteration over simple
vectors, lists and ranges are done inline, behind a test of the
arguments' classes; and slots are read directly. Two things defeat the
fast paths:

\begin{itemize}
\item Each inline operation is guarded by a flag that is true only while
  the generic functions behind it have no methods but the predefined
  ones that could apply to those arguments. Adding a method to
  \code{+}, \code{element}, \code{<} and the like for a built-in class
  turns the fast path off for the rest of the session.
\item Redefining things as the program runs (adding methods, adding
  slots) invalidates dispatch caches, which are then rebuilt on the next
  calls.
\end{itemize}

Beyond that, the advice for the compiler applies here too, in smaller
measure: local exits are cheaper than escaping ones, and code that works
on known collection classes is faster. For a program whose speed matters,
develop it interactively and compile it with \cmd{dylanc}.

\section{Size}\label{sec:perf-size}

\index{size builds}
By default the library is compiled apart from the program and linked
whole (Section~\ref{sec:dylanc-cache}), which makes builds fast but
executables large: hello world is about 460~KB of code and data. With
\cmd{-Os}, the library is compiled with the program and shaken
(Section~\ref{sec:internals-shaking}), and hello world is about 220~KB,
stripped; Section~\ref{sec:dylanc-size} compares the builds. The
original goal was under 100~KB. The shaking is by name and class, so a
call of a generic function such as \code{print-object}, \code{make} or
\code{as} keeps every method that could apply to a class with instances,
and those reach much of the library; getting further needs an analysis
of which classes flow to which calls, which the compiler does not yet
do.

A program built \cmd{-Os} may also run slower than one built
\cmd{-O2}, as C optimized for size often does: the vector-summing loop
of Section~\ref{sec:perf-c} took 0.27~s built \cmd{-Os} against 0.16~s,
though \code{tak} and the float loop ran at nearly the same speed.

\section{Measuring}\label{sec:perf-measuring}

\index{perf@\cmd{perf}}\index{profiling}
An ordinary build is compiled with \cmd{-g}, so the usual tools for C
programs work on compiled Dylan, and its \code{\#line} directives make
them name Dylan source lines. On Linux, \cmd{perf} is the simplest:

\begin{shell}
$ bin/dylanc takp.dyl
$ perf stat -e task-clock ./takp
15
             78.05 msec task-clock:u
       0.078615932 seconds time elapsed
$ perf record ./takp
$ perf report --stdio
    50.04%  takp     takp                  [.] fn1190_tak
    49.68%  takp     takp                  [.] dy_gf_tak_e10c5970_dispatch
...
\end{shell}

The functions are named as the compiler names them
(Section~\ref{sec:internals-reading}): \code{fn1190\_tak} is a method
\code{tak}; \code{dy\_gf\_tak\_e10c5970\_dispatch} is the dispatch
function of the generic function \code{tak}. Here, with the method
specialized on \code{<integer>}, half the time goes to dispatch --- the
case of the second table above. Other names to recognize:
\code{\_spec} marks a specialized entry, \code{dy\_maker\_}\meta{class}
a class's maker, and \code{GC\_}\meta{name} the collector; time in the
collector's functions is time spent allocating and collecting.
\cmd{perf annotate} shows the hot instructions with the Dylan lines they
came from.

For timing inside a program, the \code{dylan} module exports
\code{(current-milliseconds)}, an extension in both implementations that
answers a millisecond count from a fixed but arbitrary moment, as the
benchmark driver uses it.

\section{Tuning the garbage collector}\label{sec:perf-gc}

\index{garbage collection!tuning}\index{Memory Pool System!tuning}
The collector is configured by environment variables, which it reads
when the program starts; nothing in Dylan needs to change. The Memory
Pool System, the default collector (Section~\ref{sec:dylanc-gc}), reads
these:

\begin{description}
\item[\code{DYLAN\_MPS\_NURSERY}] The nursery's size in kilobytes,
  131\,072 (128~MB) by default. New objects go here, and it is collected
  when full: a larger nursery collects less often, and more of what it
  holds has died by then. A much smaller one collects too often, each
  time protecting again the older objects the program is writing into
  and taking a fault on each page it then writes; with 32~MB, programs
  that build and keep writing into a large table or vector ran slower
  than under the Boehm collector, and with 128~MB faster. The nursery is
  also the floor of the memory a program that allocates heavily uses: a
  program that must run in little memory can set, say, 32\,768.
\item[\code{DYLAN\_MPS\_EXTEND}] How much memory, in bytes, a pool
  takes from the system at a time: a megabyte by default. The
  collector's own default, which is small, costs a system call for each.
\item[\code{DYLAN\_MPS\_SPARE}] The fraction of the memory the
  collector holds that it may keep unused rather than return to the
  system, between 0 and 1; the default is 1, keeping everything. Memory
  returned must be faulted in afresh when the nursery next fills, which
  for an allocation-heavy program took longer than the collecting: with
  0.75, the MPS's own default, the benchmarks of
  Section~\ref{sec:perf-benchmarks} took 15 to 20 percent longer and used
  a quarter to a half less memory. A program's memory under the MPS is otherwise
  about twice what it is under the Boehm collector, since the collector
  lets the heap double between collections of everything.
\item[\code{DYLAN\_MPS\_PAUSE}] The longest the collector should
  work at once, in seconds, before giving the program a turn; the
  collector's default is 0.1. The MPS collects incrementally, protecting
  the memory it has yet to scan; a shorter pause means more interruptions
  and more protection faults.
\end{description}

\paragraph{The Boehm-Demers-Weiser collector.} \index{Boehm collector!tuning}
A program built with \cmd{-{}-gc boehm} reads the Boehm collector's
variables instead. These were checked with that collector (version
8.2.6), with a program that builds and discards 200 lists of 100\,000
elements. \code{GC\_PRINT\_STATS=1} makes the collector report each
collection on the standard error, which is how the counts below were
made.

\begin{description}
\item[\code{GC\_INITIAL\_HEAP\_SIZE}] The heap's starting size, in bytes
  or with a suffix such as \code{64M} or \code{1G}. The heap starts small
  and grows; a larger start means fewer collections. The test program
  made 202 collections by default and 10 with \code{64M}. Fewer
  collections are not automatically faster: here the larger heap made
  the run slower (0.24~s against 0.18~s), because touching fresh memory
  costs more than the small collections it saved. It helps programs that
  keep much live data.
\item[\code{GC\_FREE\_SPACE\_DIVISOR}] How eagerly the heap grows rather
  than collecting; the default is 3, and smaller values grow the heap
  more and collect less. With 1, the test program made 101 collections
  instead of 202 and used about 13~MB instead of 7~MB.
\item[\code{GC\_MARKERS}] How many threads mark in parallel during a
  collection. The runtime sets two, since more spend longer waiting on
  each other than marking; \code{GC\_MARKERS=1} turns parallel marking
  off, and a larger number is allowed.
\item[\code{GC\_PRINT\_STATS}] Set to 1 to report each collection.
\end{description}

\begin{shell}
$ GC_PRINT_STATS=1 ./alloc 2>&1 | grep -c 'Marking for collection'
202
$ GC_INITIAL_HEAP_SIZE=64M GC_PRINT_STATS=1 ./alloc 2>&1 |
      grep -c 'Marking for collection'
10
\end{shell}

The collectors' own documentation lists further variables. As with the
compiler's options, measure your program before and after: the right
settings depend on how much it allocates and how much it keeps.

\fussy
